% =================================================
% TikZ-рисунки для листов по логарифмическому неравенству
% =================================================

\newcommand{\logchainstudentfigure}{%
  \begin{center}
    \small
    \setlength{\fboxsep}{5pt}%
    \begin{minipage}{0.97\textwidth}
      \centering
      \colorbox{PrimaryLight}{%
        \parbox{0.92\textwidth}{%
          \centering
          \(\displaystyle
            \lg^4((x^2-4)^2)-\lg^2((x^2-4)^4)\ge192
          \)
        }%
      }

      \vspace{1mm}
      {\color{Accent}\(\Downarrow\)}\quad
      \textsf{применяем формулы с модулем}
      \vspace{1mm}

      \colorbox{PrimaryLight}{%
        \parbox{0.92\textwidth}{%
          \centering
          \(\displaystyle
            \begin{gathered}
              \text{После применения формул с модулем:}\\[1mm]
              \answerblank[118mm]
            \end{gathered}
          \)
        }%
      }

      \vspace{1mm}
      {\color{Accent}\(\Downarrow\)}\quad
      \textsf{первая замена: \(y=\answerblank[34mm]\)}
      \vspace{1mm}

      \colorbox{PrimaryLight}{%
        \parbox{0.92\textwidth}{%
          \centering
          \(\displaystyle
            \begin{gathered}
              \text{Неравенство относительно }y:\\[1mm]
              \answerblank[95mm]
            \end{gathered}
          \)
        }%
      }

      \vspace{1mm}
      {\color{Accent}\(\Downarrow\)}\quad
      \textsf{вторая замена: \(t=\answerblank[18mm]\),
        \(t\answerblank[8mm]0\)}
      \vspace{1mm}

      \colorbox{PrimaryLight}{%
        \parbox{0.92\textwidth}{%
          \centering
          \(\displaystyle
            \begin{gathered}
              \text{Квадратное неравенство:}\\[1mm]
              \answerblank[65mm]
            \end{gathered}
          \)
        }%
      }

      \vspace{2mm}
      {\sffamily\color{GrayText}После решения возвращаемся в обратном порядке:}
      \[
        t\longrightarrow y\longrightarrow
        \lg|x^2-4|\longrightarrow |x^2-4|\longrightarrow x.
      \]
    \end{minipage}
  \end{center}
}

\newcommand{\logchainteacherfigure}{%
  \begin{center}
    \small
    \setlength{\fboxsep}{5pt}%
    \begin{minipage}{0.97\textwidth}
      \centering
      \colorbox{PrimaryLight}{\parbox{0.92\textwidth}{\centering
        \(\displaystyle
          \lg^4((x^2-4)^2)-\lg^2((x^2-4)^4)\ge192
        \)}}

      \vspace{1mm}
      {\color{Accent}\(\Downarrow\)}\quad
      \textsf{применяем формулы с модулем}
      \vspace{1mm}

      \colorbox{PrimaryLight}{\parbox{0.92\textwidth}{\centering
        \(\displaystyle
          \bigl(2\lg|x^2-4|\bigr)^4
          -\bigl(4\lg|x^2-4|\bigr)^2\ge192
        \)}}

      \vspace{1mm}
      {\color{Accent}\(\Downarrow\)}\quad
      \textsf{первая замена: \(y=\lg|x^2-4|\)}
      \vspace{1mm}

      \colorbox{PrimaryLight}{\parbox{0.92\textwidth}{\centering
        \(\displaystyle
          (2y)^4-(4y)^2\ge192
          \quad\Longleftrightarrow\quad y^4-y^2-12\ge0
        \)}}

      \vspace{1mm}
      {\color{Accent}\(\Downarrow\)}\quad
      \textsf{вторая замена: \(t=y^2\), \(t\ge0\)}
      \vspace{1mm}

      \colorbox{PrimaryLight}{\parbox{0.92\textwidth}{\centering
        \(\displaystyle t^2-t-12\ge0\)}}

      \vspace{2mm}
      {\sffamily\color{GrayText}После решения возвращаемся в обратном порядке:}
      \[
        t\longrightarrow y\longrightarrow
        \lg|x^2-4|\longrightarrow |x^2-4|\longrightarrow x.
      \]
    \end{minipage}
  \end{center}
}

\newcommand{\notationtrianglefigure}{%
  \begin{center}
    \begin{tikzpicture}[node distance=7mm and 9mm]
      \node[draw=Primary,fill=PrimaryLight,rounded corners=2mm,
        text width=42mm,minimum height=16mm,align=center,font=\small\sffamily] (a)
        {$\lg^2 A=(\lg A)^2$\\степень значения функции};
      \node[draw=Secondary,fill=SecondaryLight,rounded corners=2mm,
        text width=42mm,minimum height=16mm,align=center,font=\small\sffamily,right=of a] (b)
        {$\lg(A^2)$\\степень аргумента};
      \node[draw=Accent,fill=AccentLight,rounded corners=2mm,
        text width=42mm,minimum height=16mm,align=center,font=\small\sffamily,right=of b] (c)
        {$\log_2 A$\\основание логарифма равно $2$};
    \end{tikzpicture}
  \end{center}
}

\newcommand{\argregimesstudentfigure}{%
  \begin{center}
    \begin{tikzpicture}[x=1cm,y=1cm]
      \draw[-{Stealth[length=2.5mm]},draw=GrayText] (0,0)--(11.4,0) node[right] {$B=|x^2-4|$};
      \foreach \x in {0.7,3.6,8.7}{\draw[GrayText] (\x,-0.14)--(\x,0.14);}
      \node[below=4pt,font=\small] at (0.7,0) {$0$};
      \node[below=4pt,font=\small] at (3.6,0) {$10^{-2}$};
      \node[below=4pt,font=\small] at (8.7,0) {$10^2$};
      \node[above=7pt,font=\small\sffamily,text=GrayText] at (2.1,0)
        {малые положительные значения};
      \node[above=7pt,font=\small\sffamily,text=GrayText] at (9.8,0)
        {большие значения};
    \end{tikzpicture}
  \end{center}
}

\newcommand{\argregimesteacherfigure}{%
  \begin{center}
    \begin{tikzpicture}[x=1cm,y=1cm]
      \draw[-{Stealth[length=2.5mm]},draw=GrayText] (0,0)--(11.4,0) node[right] {$B=|x^2-4|$};
      \draw[Secondary,line width=2.5pt] (0.72,0)--(3.6,0);
      \draw[Secondary,line width=2.5pt] (8.7,0)--(11.15,0);
      \node[circle,draw=Accent,fill=white,line width=0.9pt,inner sep=2pt] at (0.7,0) {};
      \node[circle,draw=Secondary,fill=Secondary,inner sep=2pt] at (3.6,0) {};
      \node[circle,draw=Secondary,fill=Secondary,inner sep=2pt] at (8.7,0) {};
      \foreach \x in {0.7,3.6,8.7}{\draw[GrayText] (\x,-0.14)--(\x,0.14);}
      \node[below=4pt,font=\small] at (0.7,0) {$0$};
      \node[below=4pt,font=\small] at (3.6,0) {$10^{-2}$};
      \node[below=4pt,font=\small] at (8.7,0) {$10^2$};
      \node[above=7pt,font=\small\bfseries\sffamily,text=Secondary] at (2.1,0)
        {$0<B\le10^{-2}$};
      \node[above=7pt,font=\small\bfseries\sffamily,text=Secondary] at (9.8,0)
        {$B\ge10^2$};
    \end{tikzpicture}
  \end{center}
}

\newcommand{\finalnumberlinestudentfigure}{%
  \begin{center}
    \begin{tikzpicture}[x=0.72cm,y=1cm]
      \draw[-{Stealth[length=2.5mm]},draw=GrayText] (-10.5,0)--(10.5,0) node[right] {$x$};
      \foreach \x in {-8,-5.7,-3.7,-1.8,1.8,3.7,5.7,8}{\draw[GrayText] (\x,-0.15)--(\x,0.15);}
      \node[below=5pt,font=\scriptsize] at (-8,0) {$-2\sqrt{26}$};
      \node[below=16pt,font=\scriptsize] at (-5.7,0) {$-\frac{\sqrt{401}}{10}$};
      \node[below=5pt,font=\scriptsize] at (-3.7,0) {$-2$};
      \node[below=16pt,font=\scriptsize] at (-1.8,0) {$-\frac{\sqrt{399}}{10}$};
      \node[below=16pt,font=\scriptsize] at (1.8,0) {$\frac{\sqrt{399}}{10}$};
      \node[below=5pt,font=\scriptsize] at (3.7,0) {$2$};
      \node[below=16pt,font=\scriptsize] at (5.7,0) {$\frac{\sqrt{401}}{10}$};
      \node[below=5pt,font=\scriptsize] at (8,0) {$2\sqrt{26}$};
      \node[above=5pt,font=\footnotesize\sffamily,text=GrayText] at (0,0.5)
        {Отметьте открытые и закрытые границы, затем выделите решение.};
    \end{tikzpicture}
  \end{center}
}

\newcommand{\finalnumberlineteacherfigure}{%
  \begin{center}
    \begin{tikzpicture}[x=0.72cm,y=1cm]
      \draw[-{Stealth[length=2.5mm]},draw=GrayText] (-10.5,0)--(10.5,0) node[right] {$x$};
      \draw[Secondary,line width=2.5pt] (-10.2,0)--(-8,0);
      \draw[Secondary,line width=2.5pt] (-5.7,0)--(-3.73,0);
      \draw[Secondary,line width=2.5pt] (-3.67,0)--(-1.8,0);
      \draw[Secondary,line width=2.5pt] (1.8,0)--(3.67,0);
      \draw[Secondary,line width=2.5pt] (3.73,0)--(5.7,0);
      \draw[Secondary,line width=2.5pt] (8,0)--(10.2,0);
      \node[circle,fill=Secondary,draw=Secondary,inner sep=2.1pt] at (-8,0) {};
      \node[circle,fill=Secondary,draw=Secondary,inner sep=2.1pt] at (-5.7,0) {};
      \node[circle,draw=Accent,fill=white,line width=0.9pt,inner sep=2.1pt] at (-3.7,0) {};
      \node[circle,fill=Secondary,draw=Secondary,inner sep=2.1pt] at (-1.8,0) {};
      \node[circle,fill=Secondary,draw=Secondary,inner sep=2.1pt] at (1.8,0) {};
      \node[circle,draw=Accent,fill=white,line width=0.9pt,inner sep=2.1pt] at (3.7,0) {};
      \node[circle,fill=Secondary,draw=Secondary,inner sep=2.1pt] at (5.7,0) {};
      \node[circle,fill=Secondary,draw=Secondary,inner sep=2.1pt] at (8,0) {};
      \foreach \x in {-8,-5.7,-3.7,-1.8,1.8,3.7,5.7,8}{\draw[GrayText] (\x,-0.15)--(\x,0.15);}
      \node[below=5pt,font=\scriptsize] at (-8,0) {$-2\sqrt{26}$};
      \node[below=16pt,font=\scriptsize] at (-5.7,0) {$-\frac{\sqrt{401}}{10}$};
      \node[below=5pt,font=\scriptsize] at (-3.7,0) {$-2$};
      \node[below=16pt,font=\scriptsize] at (-1.8,0) {$-\frac{\sqrt{399}}{10}$};
      \node[below=16pt,font=\scriptsize] at (1.8,0) {$\frac{\sqrt{399}}{10}$};
      \node[below=5pt,font=\scriptsize] at (3.7,0) {$2$};
      \node[below=16pt,font=\scriptsize] at (5.7,0) {$\frac{\sqrt{401}}{10}$};
      \node[below=5pt,font=\scriptsize] at (8,0) {$2\sqrt{26}$};
    \end{tikzpicture}
  \end{center}
}
